\documentclass[10pt,a4paper]{amsart}
\usepackage{amsmath,amssymb,amsthm,graphicx}
% Retained minimal preamble for the Stacks Project chapter sites-cohomology.tex (Cohomology on Sites).
% Verbatim excerpt of the chapter preamble: Xy-pic setup (native preamble lines 14-18) and the native
% theorem declarations with the mathematical macros (native preamble lines 154-217). Copied mechanically;
% nothing added, removed or redefined. The wrapper class and the amsmath/amssymb/amsthm/graphicx base
% packages are supplied by the assembler, so the class/externaldocument part of the native preamble is omitted.
% For commutative diagrams we use Xy-pic
\usepackage[all]{xy}

% We use 2cell for 2-commutative diagrams.
\xyoption{2cell}
% Theorem environments.
%
\theoremstyle{plain}
\newtheorem{theorem}[subsection]{Theorem}
\newtheorem{proposition}[subsection]{Proposition}
\newtheorem{lemma}[subsection]{Lemma}

\theoremstyle{definition}
\newtheorem{definition}[subsection]{Definition}
\newtheorem{example}[subsection]{Example}
\newtheorem{exercise}[subsection]{Exercise}
\newtheorem{situation}[subsection]{Situation}

\theoremstyle{remark}
\newtheorem{remark}[subsection]{Remark}
\newtheorem{remarks}[subsection]{Remarks}

\numberwithin{equation}{subsection}

% Macros
%
\def\lim{\mathop{\mathrm{lim}}\nolimits}
\def\colim{\mathop{\mathrm{colim}}\nolimits}
\def\Spec{\mathop{\mathrm{Spec}}}
\def\Hom{\mathop{\mathrm{Hom}}\nolimits}
\def\Ext{\mathop{\mathrm{Ext}}\nolimits}
\def\SheafHom{\mathop{\mathcal{H}\!\mathit{om}}\nolimits}
\def\SheafExt{\mathop{\mathcal{E}\!\mathit{xt}}\nolimits}
\def\Sch{\mathit{Sch}}
\def\Mor{\mathop{\mathrm{Mor}}\nolimits}
\def\Ob{\mathop{\mathrm{Ob}}\nolimits}
\def\Sh{\mathop{\mathit{Sh}}\nolimits}
\def\NL{\mathop{N\!L}\nolimits}
\def\CH{\mathop{\mathrm{CH}}\nolimits}
\def\proetale{{pro\text{-}\acute{e}tale}}
\def\etale{{\acute{e}tale}}
\def\QCoh{\mathit{QCoh}}
\def\Ker{\mathop{\mathrm{Ker}}}
\def\Im{\mathop{\mathrm{Im}}}
\def\Coker{\mathop{\mathrm{Coker}}}
\def\Coim{\mathop{\mathrm{Coim}}}

% Boxtimes
%
\DeclareMathSymbol{\boxtimes}{\mathbin}{AMSa}{"02}

%
% Macros for moduli stacks/spaces
%
\def\QCohstack{\mathcal{QC}\!\mathit{oh}}
\def\Cohstack{\mathcal{C}\!\mathit{oh}}
\def\Spacesstack{\mathcal{S}\!\mathit{paces}}
\def\Quotfunctor{\mathrm{Quot}}
\def\Hilbfunctor{\mathrm{Hilb}}
\def\Curvesstack{\mathcal{C}\!\mathit{urves}}
\def\Polarizedstack{\mathcal{P}\!\mathit{olarized}}
\def\Complexesstack{\mathcal{C}\!\mathit{omplexes}}
% \Pic is the operator that assigns to X its picard group, usage \Pic(X)
% \Picardstack_{X/B} denotes the Picard stack of X over B
% \Picardfunctor_{X/B} denotes the Picard functor of X over B
\def\Pic{\mathop{\mathrm{Pic}}\nolimits}
\def\Picardstack{\mathcal{P}\!\mathit{ic}}
\def\Picardfunctor{\mathrm{Pic}}
\def\Deformationcategory{\mathcal{D}\!\mathit{ef}}
\usepackage{cleveref}
\crefname{originaltheorem}{original theorem}{original theorems}
\Crefname{originaltheorem}{Original theorem}{Original theorems}
\makeatletter
\newlabel{topology-theorem-characterize-proper}{{005R}{source}{}{sourcepointer-1}{}}
\newlabel{topology-theorem-characterize-proper@cref}{{[originaltheorem][0][]005R}{[1][0][]source}{}{}{}}
\makeatother
\title{A proper surjective morphism of locally quasi-compact spaces is a qc covering}
\author{The Stacks Project Authors}
\date{stacks-project:sites-cohomology.tex:local-snapshot-20261011; GNU FDL 1.2 or later}
\begin{document}
\maketitle
\noindent\textit{Editorial note.} All delivered mathematical text is the author's own, copied from the pinned Stacks Project chapter file sites-cohomology.tex; no statement or proof step was rewritten, reordered or removed.\par\medskip
\noindent\textit{Editorial note.} The proof is the author's complete argument: the finite covering of the fibre, the closed set Z and the final inclusion are all inside the counted span.\par\medskip
\noindent\textit{Editorial note.} The label lemma-proper-surjective-is-qc-covering occurs in no other staged Stacks chapter, so no cross-chapter restatement had to be weighed for this card.\par\medskip
\noindent\textit{Editorial note.} The single labelled reference occurring inside the counted spans points to a chapter of the Project not staged in this card's source tree; it is carried in assembly.external\_references as a non-load-bearing navigation pointer recorded at the invoking line.\par\medskip
\noindent\textit{Editorial note.} The preamble is the chapter's own retained minimal preamble (sources/stacks-sites-cohomology/retained-preamble.tex), a verbatim excerpt of the native preamble: the Xy-pic setup and the native theorem declarations with the mathematical macros.\par\medskip
\noindent\textit{Allowed standard prerequisites (editorial).}\begin{itemize}
\item The category LC of Hausdorff locally quasi-compact spaces and its quasi-compact subsets. Same chapter, the category LC and the notion of a quasi-compact subset used in the definition of qc coverings. Supplies the ambient category and the quasi-compact neighbourhoods chosen in the proof.
\item qc covering. Same chapter, the definition of a qc covering as a family such that every point has a neighbourhood covered by finitely many quasi-compact subsets of the sources. Supplies the property verified for the single morphism f.
\item Properness: closedness and quasi-compactness of geometric fibres. Topology, the characterisation of proper maps by closedness together with quasi-compactness of the fibres. Supplies both the closedness of f and the quasi-compactness of the fibre X\_y that the proof uses.
\end{itemize}\medskip
\ifdefined\hypertarget\hypertarget{sourcepointer-1}{}\fi
\noindent\textit{Source reference.} Original theorem 005R: sources/stacks-sites-cohomology/pinned-original-source.tex, source lines 7599--7599. Cross-chapter navigation pointer invoked at proof line 7599, where the author uses that a proper map is closed. The target chapter is not staged in this card's source tree, so the pointer records the invoking line; it only resolves the reference for the reader, and no bibliography citation is involved.\par\medskip
\begin{lemma}
\label{lemma-proper-surjective-is-qc-covering}
Let $f : X \to Y$ be a morphism of $\textit{LC}$.
If $f$ is proper and surjective, then $\{f : X \to Y\}$
is a qc covering.
\end{lemma}

\begin{proof}
Let $y \in Y$ be a point. For each $x \in X_y$ choose a quasi-compact
neighbourhood $E_x \subset X$. Choose $x \in U_x \subset E_x$ open.
Since $f$ is proper the fibre $X_y$ is quasi-compact and we find
$x_1, \ldots, x_n \in X_y$ such that
$X_y \subset U_{x_1} \cup \ldots \cup U_{x_n}$.
We claim that $f(E_{x_1}) \cup \ldots \cup f(E_{x_n})$ is a neighbourhood of
$y$. Namely, as $f$ is closed
(Topology, Theorem \ref{topology-theorem-characterize-proper})
we see that $Z = f(X \setminus U_{x_1} \cup \ldots \cup U_{x_n})$
is a closed subset of $Y$ not containing $y$. As $f$ is surjective
we see that $Y \setminus Z$ is contained in
$f(E_{x_1}) \cup \ldots \cup f(E_{x_n})$ as desired.
\end{proof}

\end{document}
